\begin{abstract}
Standard foundation model training relies on a rigid sequential paradigm: Next-Token Prediction (NTP) pre-training, followed by Supervised Fine-Tuning (SFT), and finally Reinforcement Learning (RL). While recent work demonstrates that RL excursions during pre-training can expand model distribution coverage, fixed or manual excursion schedules fail to account for dynamic checkpoint readiness. In this paper, we investigate the underlying mechanisms of reinforcement learning plasticity across intermediate pre-training checkpoints. We establish empirically across verified open model training runs that pre-RL state signals—specifically gradient alignment between NTP loss and task rewards ($\cos(\mathbf{g}_{\text{NTP}}, \mathbf{g}_{\text{RL}})$, Pearson $r=0.7562, p=1.92\times 10^{-5}$, Spearman $\rho=0.8718, p=2.89\times 10^{-8}$), policy entropy, and baseline accuracy—predict subsequent RL performance gains before compute is expended. We show that a predictive linear model trained on one open model family (\texttt{SmolLM-135M}) transfers zero-shot to predict RL plasticity on an unseen model architecture (\texttt{distilgpt2}, zero-shot $R^2=+0.2738$, Spearman $\rho=+0.9519$). All empirical records, parameter hashes, and raw generation logs are validated for provenance under \texttt{artifacts/empirical/}.
\end{abstract}
\keywords{reinforcement learning, large language models, training dynamics, optimization, compute allocation, empirical provenance}
